\section{Evaluation design}
\label{sec:evaluation}
We test whether the shared core learns predictive and semantic features,
how much mutable state online adaptation requires, and which component changes
behavior. The studies use different experimental units, listed with their
released records in Appendix~\ref{app:protocols}.

\paragraph{Representation and joint learning.}
The sensory sources include SHD/SSC audio \citep{cramer2022heidelberg}, retinal
digit images and rendered words, and a causal video-feature cache. Action
sequences use Meta-World reach, push, pick-place, and door-open tasks
\citep{yu2019metaworld}. Five action initializations sharing one sensory
foundation compare sequential training, action-only updating, and joint
training. Three additional controlled event-model studies examine temporal
structure, prediction of future observation statistics, and image/audio factor
association. They use separate probe models and three seeds each. These probe models are separate from the checkpoints used in the deployment
comparison.

\paragraph{The main deployment stream.}
Three independently trained foundations, each with three action seeds, produce
nine paired runs. Each uses the sequential checkpoint. All methods see the
same six-stage layout: source, attenuation, source, shear, attenuation, and
source return. Every stage includes the four manipulation tasks. The tested
shift values belong to the registered bank. Shifted trajectories come from an
execution archive separate from router calibration. Source stages revisit
calibration episodes, so source-return measurements quantify calibrated reentry.

The five configurations are frozen prediction, external state, action-adapter
tuning, prediction-path tuning, and prediction-path tuning with replay. All
three trainable comparison arms use Adam. ``Prediction-path tuning'' refers to
the deployed prediction path, not every parameter in the multimodal model.
We count mutable parameters, persistent optimizer tensors, state and routing
buffers, and replay where applicable. We also record update latency using
CUDA-synchronized measurements after warm-up.

The main study replays a common recorded trajectory in causal order: each
prediction is evaluated before its outcome is used to update the method.
Each run contains 348 transitions, of which 174 are shifted. The primary
score uses the 126 shifted transitions remaining after excluding the first
four transitions of every episode, which establish its route.
The score is the relative reduction in prediction mean squared error
(MSE). If $E_{\mathrm{fr}}$ and $E_M$ denote frozen and
method-$M$ errors on the same evaluated transitions, prediction gain is
\begin{equation}
 G_M=100\left(1-\frac{E_M}{E_{\mathrm{fr}}}\right)\%.
 \label{eq:gain}
\end{equation}
A positive gain means smaller error. Action tracking applies the known actuator
map to each corrected command and measures its distance from the desired
command. This replay metric is separate from reward on method-specific live
trajectories, evaluated below. Source-return correction measures
the size of unnecessary command adjustments upon calibrated reentry. Action
seeds are averaged within each foundation. Confidence intervals (CIs) use
20,000 bootstrap draws over the three foundation means, preserving paired
method comparisons.

\paragraph{Component and closed-loop tests.}
The component study compares full state, score-only state, and removal of the
inverse. A separate 36-trace check verifies that bounding the routing history
preserves the original implementation's observable behavior. Closed-loop
attenuation uses five model seeds, three tasks, and five environment replicas:
75 trajectories per method. The tasks are reach, push, and pick-place, with
attenuation gain 0.4. Each trajectory contains 13 decisions over 104 environment
steps. Adaptation state carries across the five replicas within each model/task
pair; a new controller is created for the next pair. Reward and success are
evaluation outputs, never adaptation inputs. Its CIs use paired resampling
of model, task, and replica axes. The live-control tuning and replay controllers
are defined separately in Appendix~\ref{app:control}.

\paragraph{Classical calibration and training-stage diagnostics.}
A matched six-stage replay compares the external state with routed actuator-space
recursive least squares (RLS), using a two-by-two estimate for action coupling.
The closed-loop study includes a separate scalar RLS reference for attenuation.
Appendix~\ref{app:additional} also reports a one-foundation development study
of direct output-space calibration, including normalized least mean squares
(NLMS), RLS, and a low-rank readout update. We report each comparison under its own protocol because the calibrators
correct different quantities.
A further one-foundation pilot applies the same external-state interface to
different offline training stages.
